\section{Supplementary material: proof manual}\label{sec:appendix}

The supplementary material is a complete proof manual, targeted at
100--150 pages, and is organised as follows.

\subsection{Page budget}

\begin{center}
\begin{tabular}{ll}
\toprule
Chapter & Planned pages \\
\midrule
String-flow vocabulary & 15--20 \\
Block calculus & 20--30 \\
Unified core proof & 30--40 \\
Divergence bridge & 15--20 \\
Lean theorem index & 15--25 \\
Dependency graph and ledger & 5--10 \\
\midrule
Total & 100--145 \\
\bottomrule
\end{tabular}
\end{center}

The current compiled companion draft is
\texttt{paper/supplement\_manual.tex}; it currently contains the
vocabulary, block calculus, unified core, bridge, Lean index,
dependency ledger, finite base, CRT layer, and interface chapters,
and is maintained incrementally toward the target length.

\subsection{Chapter contents and current status}

\begin{center}
\small
\begin{tabular}{@{}p{0.24\textwidth}p{0.48\textwidth}p{0.22\textwidth}@{}}
\toprule
Chapter & Contents & Status \\
\midrule
1 & vocabulary, PMI, block equations & draft \\
2 & unified core, segment exclusions, CRT & draft \\
3 & cycle bridge, corrected window, failure window & draft \\
4 & Lean theorem index & partial list \\
5 & dependency graph and ledger & draft \\
6 & finite-prefix base and explicit expansion & draft \\
7 & CRT candidate and terminal residue & draft \\
8 & Lean bridge interfaces & draft \\
9 & block calculus and reset equations & draft \\
10 & PMI and PMI-B budgets & draft \\
11 & arithmetic counterexample and corrected window & draft \\
12 & failure window construction & draft \\
13 & worked computations ledger & draft \\
14 & theorem-by-theorem status matrix & draft \\
15 & genericity ledger for $(a,b)$ & draft \\
16 & cycle word algebra & draft \\
17 & corrected window bound architecture & draft \\
18 & Lean build and reproducibility guide & draft \\
19 & block-head candidate exclusion and remaining bridge & draft \\
20 & finite-prefix verification details & draft \\
21 & divergence bridge in full & draft \\
22 & rise blocks, C3 blocks, and block-layer balance & draft \\
23 & glossary, notation, and reading guide & draft \\
24 & proof architecture and dependency diagrams & draft \\
25 & why the proof is exact & draft \\
26 & full derivation of the $d=2$ exclusion & draft \\
27 & full derivation of the $d=3$ exclusion & draft \\
28 & full derivation of the $d\ge4$ exclusions & draft \\
29 & finite-prefix proof script outline & draft \\
30 & CRT layer: statement-by-statement commentary & draft \\
31 & corrected reachability predicate & draft \\
32 & worked examples of the exact ledger & draft \\
33 & submission and release checklist & draft \\
\bottomrule
\end{tabular}
\end{center}

\subsection{Reading guide}

\begin{enumerate}[leftmargin=*]
\item Readers who want the exact vocabulary should begin with the
  string-flow chapter and check the Lean correspondence table in
  \S2 of the main paper.
\item Readers who want the proof of the unified core should begin with
  the core chapter, then follow the segment-exclusion details.
\item Readers who want the formalisation status should begin with the
  Lean index and the dependency graph.
\item The status ledger is authoritative: a lemma marked pending must
  not be cited as closed.
\end{enumerate}

\subsection{Manual outline}

\begin{enumerate}[leftmargin=*]
\item \textbf{String-flow vocabulary.}  Exact definitions of words,
validity, weight, molecules, prefix margins, and the PMI identity.
\item \textbf{Block calculus.}  Rise blocks, C3 blocks, reset equations,
block-head candidates, and the corrected reachability predicate.
\item \textbf{Unified core proof.}  Full derivations for
\S36.30.22 and \S36.30.23, including the candidate parameterisation,
segment-length exclusions, and the finite-prefix base.
\item \textbf{Divergence bridge.}  The corrected window bound, failure
windows, and the no-cycle assembly.
\item \textbf{Lean index.}  A theorem-by-theorem index of the files
\texttt{FinitePrefix.lean}, \texttt{UnifiedCoreBridge.lean},
\texttt{UnifiedCoreAudit.lean}, \texttt{CycleBridge.lean}, and
\texttt{DwdbDiv.lean}.
\item \textbf{Dependency graph.}  Each mathematical lemma is linked to
its Lean statement and to the sections of this paper.
\end{enumerate}

\subsection{Cross-reference policy}

The supplementary material does not paste Lean files line by line.  For
each definition and theorem it gives:
\begin{itemize}[leftmargin=*]
\item the mathematical statement;
\item the Lean name;
\item the file and line region;
\item the proof status;
\item the list of dependencies.
\end{itemize}

\subsection{Reproducing the paper}

The main document is built with:
\begin{verbatim}
pdflatex main
bibtex main
pdflatex main
pdflatex main
\end{verbatim}
The supplementary manual is built with:
\begin{verbatim}
pdflatex supplement_manual
pdflatex supplement_manual
\end{verbatim}

\subsection{Manual maintenance checklist}

Whenever a Lean statement changes status, the paper and the manual are
updated together:
\begin{itemize}[leftmargin=*]
\item re-run the relevant \texttt{lake build} target;
\item re-run \texttt{\#print axioms} for changed bridge theorems;
\item update every status table and the dependency-edge table;
\item update \texttt{STATUS.md} and the page counts;
\item do not promote \texttt{unified\_core\_final\_no\_hge} or
  \texttt{FinalStatement.five\_\allowbreak x\_\allowbreak plus\_\allowbreak
  one\_\allowbreak diverges\_\allowbreak at\_\allowbreak 7} until
  their ledger entries are changed;
\item keep every mathematical statement aligned with the Lean name
  that is cited for it.
\end{itemize}

\subsection{Status ledger}

The status ledger is kept in \texttt{docs/lean\_assembly\_plan.md},
\texttt{docs/cycle\_bridge\_window\_to\_no\_cycle.md}, and
\texttt{docs/dwdb\_div\_assembly.md}.  The paper and the supplementary manual
must not claim a Lean statement is closed when its ledger entry is
pending.

\subsection{Requirements checklist}

The following table maps the manuscript requirements to their
locations and current status.

\begin{center}
\small
\begin{tabular}{@{}p{0.34\textwidth}p{0.34\textwidth}p{0.24\textwidth}@{}}
\toprule
Requirement & Location & Status \\
\midrule
project structure and \texttt{refs.bib} & \texttt{main.tex},
  \texttt{sections/*.tex} & done \\
core definitions and main theorem & \texttt{intro.tex},
  \texttt{framework.tex} & done \\
string-flow framework and PMI/PMI-B & \texttt{framework.tex} & done \\
full vs general orbit reachability & \texttt{framework.tex} & done \\
unified core outline and statuses & \texttt{core.tex} & done \\
divergence bridge and conditional chain & \texttt{bridge.tex} & done \\
Lean verification setup and audit & \texttt{lean.tex} & done \\
supplementary proof manual & \texttt{supplement\_manual.tex} &
  102 pages \\
constants $13$ and thresholds $+10,+11,+12$ & throughout & preserved \\
forbidden phrases absent & throughout & verified \\
pending statements not promoted & throughout & verified \\
Lean build final confirmation & Lean repository & pending \\
\bottomrule
\end{tabular}
\end{center}

The final row is the release gate: the manuscript is complete as a
draft with the pending statuses recorded, but the final authority is
the Lean build, which is not yet closed.
